% ===== PRÉAMBULE CANONIQUE — à coller VERBATIM en tête de CHAQUE .tex =====
\documentclass[11pt,a4paper]{article}
\usepackage[utf8]{inputenc}\usepackage[T1]{fontenc}\usepackage{lmodern}
\usepackage[french]{babel}
\usepackage{amsmath,amssymb,mathtools}\usepackage{icomma}
\usepackage[version=4]{mhchem}          % \ce{...} : équations & formules
\usepackage{siunitx}\sisetup{locale=FR} % \qty{}{}, \si{}, \ang{}
\usepackage{chemfig}                    % \chemfig{...} : structures développées
\usepackage{xcolor}\usepackage{colortbl}
\usepackage{tikz}\usepackage{pgfplots}\pgfplotsset{compat=1.18}
\usepackage[most]{tcolorbox}\usepackage{enumitem}\usepackage{array,booktabs}
\usepackage[a4paper,margin=2.2cm]{geometry}
\usetikzlibrary{arrows.meta,calc,angles,quotes,positioning,patterns,backgrounds,shapes.geometric}
\definecolor{primary}{RGB}{222,49,99}\definecolor{primarydark}{RGB}{150,22,55}
\definecolor{defcol}{RGB}{0,114,178}\definecolor{methcol}{RGB}{52,92,148}
\definecolor{excol}{RGB}{0,135,90}\definecolor{attncol}{RGB}{200,40,55}
\tcbset{boxbase/.style={enhanced,breakable,boxrule=0.9pt,arc=2.5pt,
  left=3mm,right=3mm,top=2.3mm,bottom=2.3mm,fonttitle=\bfseries,coltitle=white}}
% --- boîtes TUTORIEL (noms EXACTS, ne pas renommer) ---
\newtcolorbox{cle}[1][]{boxbase,colback=defcol!6,colframe=defcol,
  title={\textsc{À retenir}\ifx\relax#1\relax\relax\else\textmd{ : \itshape #1}\fi}}
\newtcolorbox{methode}[1][]{boxbase,colback=methcol!7,colframe=methcol,sharp corners,
  title={\textsc{Méthode}\ifx\relax#1\relax\relax\else\textmd{ --- \itshape #1}\fi}}
\newtcolorbox{exemplebox}[1][]{boxbase,colback=excol!5,colframe=excol,
  title={\textsc{Exemple}\ifx\relax#1\relax\relax\else\textmd{ : \itshape #1}\fi}}
\newtcolorbox{piege}[1][]{boxbase,colback=attncol!6,colframe=attncol,
  title={\textsc{Piège}\ifx\relax#1\relax\relax\else\textmd{ : \itshape #1}\fi}}
\newtcolorbox{remarque}[1][]{boxbase,colback=black!4,colframe=black!45,
  title={\textsc{Remarque}\ifx\relax#1\relax\relax\else\textmd{ : \itshape #1}\fi}}
% --- boîtes EXERCICE / CORRIGÉ (noms EXACTS, auto-comptées) ---
\newtcolorbox[auto counter]{exo}[1][]{boxbase,colback=primary!4,colframe=primary,
  title={\textsc{Exercice}~\thetcbcounter\ifx\relax#1\relax\relax\else\textmd{ --- \itshape #1}\fi}}
\newtcolorbox[auto counter]{corrige}[1][]{boxbase,colback=excol!4,colframe=excol,
  title={\textsc{Corrigé}~\thetcbcounter\ifx\relax#1\relax\relax\else\textmd{ --- \itshape #1}\fi}}
\newcommand{\R}{\mathbb{R}}\newcommand{\N}{\mathbb{N}}\newcommand{\Z}{\mathbb{Z}}\newcommand{\ds}{\displaystyle}
% (un \usepackage{fancyhdr}+en-tête est possible ; il est ignoré côté web)
% =========================================================================
\begin{document}

\begin{corrige}[loi de Boyle-Mariotte ($T$ constante)]
On isole l'inconnue dans $p_1 V_1 = p_2 V_2$.
\begin{enumerate}[label=\alph*)]
  \item $V_2=\dfrac{p_1 V_1}{p_2}=\dfrac{1\times6}{2}=\qty{3}{\litre}$.
  \item $p_2=\dfrac{p_1 V_1}{V_2}=\dfrac{4\times3}{6}=\qty{2}{atm}$.
  \item $p$ et $V$ inversement proportionnels : $V$ est \textbf{divisé par $3$}.
\end{enumerate}
\end{corrige}

\begin{corrige}[loi de Charles--Gay-Lussac ($p$ constante)]
On isole $V_2=V_1\times\dfrac{T_2}{T_1}$, après conversion en kelvins.
\begin{enumerate}[label=\alph*)]
  \item $V_2=2\times\dfrac{600}{300}=\qty{4}{\litre}$.
  \item $T_1=300$ K, $T_2=600$ K : $V_2=5\times\dfrac{600}{300}=\qty{10}{\litre}$.
  \item $V$ et $T$ directement proportionnels : $V$ est \textbf{doublé} ($\times2$).
\end{enumerate}
\end{corrige}

\begin{corrige}[transformation isochore ($V$ constant)]
On isole $p_2=p_1\times\dfrac{T_2}{T_1}$, températures en kelvins.
\begin{enumerate}[label=\alph*)]
  \item $p_2=1\times\dfrac{450}{300}=\qty{1.5}{atm}$.
  \item $T_1=300$ K, $T_2=150$ K : $p_2=2\times\dfrac{150}{300}=\qty{1}{atm}$.
\end{enumerate}
\end{corrige}

\begin{corrige}[loi d'Avogadro ($p$ et $T$ constants)]
On isole $V_2=V_1\times\dfrac{n_2}{n_1}$.
\begin{enumerate}[label=\alph*)]
  \item $V_2=12\times\dfrac{1{,}0}{0{,}5}=\qty{24}{\litre}$.
  \item $V$ et $n$ directement proportionnels : $V$ est \textbf{doublé} ($\times2$).
\end{enumerate}
\end{corrige}

\end{document}
